\cotefile{77}
\scanpage{129}

% L1
donc si $\unsure{Q'|q'}$ est lisse, $\delta'(q)$ définit un iso

% L2
\[
 L^{{\otimes}n}\overset{\sim}{\longleftarrow}(\det E)^{\otimes 2}.\tag*{(1.5)}
\]

% L3
Pour la suite on suppose que \unsure{$Q$ est|c'est|$q$ est} lisse

% L4
\uline{2} Groupe orthogonal\inserted{ et spécial-orthogonal}, cas $n$ impair $=2\nu+1$. // On a

% L5–6
\[
 \left\{\begin{aligned}
 q_n(x_1,y_1,\mathbin{-}x_\nu,y_\nu,z)&=\sum x_i y_i+z^2\\
 \delta'(q_n)&=(-1)^\nu \qquad \delta'(-q_n)=-\delta'(q_n)=(-1)^{\nu+1}
 \end{aligned}\right.
\]

% L6, left
On pose

% L7
\[
 O(q)=\underline{\mathrm{Aut}}(q)
\]

% L8–9
\diagram{\begin{tikzcd}[ampersand replacement=\&,column sep=small,row sep=small]
 {SO(q)=\underline{\mathrm{Aut}}(q)\cap\mathrm{Sl}(E)=\mathrm{Ker}\bigl(}
 \& {O(q)}\arrow[r,"\det"]\arrow[dr,bend right=20] \& {\mathbb G_m\bigl)}\\
 \& \& \mu_2\arrow[u,hook]
\end{tikzcd}}

% L10
\uline{NB} Sans condition sur parité de $n$, on a

% L11, left margin
\margin{\unsure{remonté|remonte|remarque}}

% L11
\[
 \mathbb G_m\cap O(q)=\mu_2\subset O(q)\qquad \mu_2=\underline{\mathrm{Centre}}\,O(q)\tag*{(2.1)}
\]

% L12
Dans le cas actuel $n$ impair, on a

% L13
\[
 \mu_2\cap SO(q)=\{1\}\,.\tag*{(2.2)}
\]

% L14–17, left rule
\uline{Thm} $O(q)$ est le produit direct des \uline{ss-groupes} $\mu_2$ et $SO(q)$, $SO(q)$ est lisse sur $S$ à fibres conn. (et en fait \deleted{simple si} réductif \deleted{si \unsure{$q$}} et $\unsure{=}\ \frac{1}{2}$ simple si $n\geq3$\inserted{ (i.e. $\nu\geq1$)}, simple si $n\geq7$\inserted{ (i.e. $\nu\geq\unsure{3}$)}).

% L18–20
Structure spéciale orthogonale sur $E$ \inserted{d'indice $\varepsilon$}, \deleted{\unsure{couple}} où $\varepsilon$ est une section de $\mathcal O_S^*$ (en pratique, $\varepsilon=\pm1$) : c'est une section-\inserted{base} $\omega$ de $\underline{\det}\,E$, telle que

% L21
\[
 \boxed{\underline{\omega}^{\otimes-2}=\varepsilon\,\delta'(q)}\tag*{(2.3)}
\]

% L21–24, right margin
\margin{Pour $n$ impair, \deleted{$q\mapsto\unsure{-q|q}$}\inserted{$(q,\omega)\mapsto(-q,\omega)$} \deleted{\unsure{fait passer}} bijection de l'ens des structures spéciales orth. pour $\varepsilon$, vers celles pour $-\varepsilon$ (donc pas g$^{\mathrm{de}}$ diff\ldots)}

% L22
On a évidemment

% L23
\[
 \mathrm{Aut}(E,q,\omega)=SO(q)\tag*{(2.4)}
\]

% L24–25
\[
 \left\{\begin{aligned}
 \text{Si}\qquad\varepsilon&=(-1)^\nu,\quad\text{alors }q_n\text{ est un prototype}\\
 \text{Si}\qquad\varepsilon&=(-1)^{\nu+1},\quad\text{---}\quad -q_n\quad\text{----------------}
 \end{aligned}\right.
\]

% lower left
% archive sheet number: 128
